\documentclass[12pt,a4paper]{article}
\usepackage[utf8]{inputenc}
\usepackage[spanish]{babel}
\usepackage{amsmath, amssymb}
\usepackage{geometry}
\usepackage{booktabs}
\usepackage{xcolor}
\usepackage{parskip}

\geometry{margin=2.5cm}

\title{\textbf{Modelado y Simulaci\'on}\\[0.5em]
       \large Tema 03 --- Interpolaci\'on de Lagrange}
\author{}
\date{}

\begin{document}
\maketitle
\thispagestyle{empty}

% ─────────────────────────────────────────────────────────────
\section*{Ejercicio 5}

Hallar el polinomio interpolante de Lagrange para:
\[
  x = [0,\;1,\;2], \qquad y = [1,\;3,\;0]
\]

% ─────────────────────────────────────────────────────────────
\section*{Resoluci\'on}

Tenemos $n+1 = 3$ puntos, por lo que el polinomio interpolante tiene grado $\leq 2$:
\[
  (x_0,\,y_0)=(0,1),\quad
  (x_1,\,y_1)=(1,3),\quad
  (x_2,\,y_2)=(2,0).
\]

\subsection*{Paso 1 \;---\; Polinomios base de Lagrange}

La f\'ormula general de cada polinomio base es:
\[
  L_i(x) = \prod_{\substack{j=0\\j\neq i}}^{n}
            \frac{x - x_j}{x_i - x_j}
\]

\textbf{$L_0(x)$} \;(ra\'ices en $x_1=1$ y $x_2=2$):
\[
  L_0(x)
    = \frac{(x-1)(x-2)}{(0-1)(0-2)}
    = \frac{(x-1)(x-2)}{(-1)(-2)}
    = \frac{(x-1)(x-2)}{2}
\]

\textbf{$L_1(x)$} \;(ra\'ices en $x_0=0$ y $x_2=2$):
\[
  L_1(x)
    = \frac{(x-0)(x-2)}{(1-0)(1-2)}
    = \frac{x(x-2)}{(1)(-1)}
    = -x(x-2)
\]

\textbf{$L_2(x)$} \;(ra\'ices en $x_0=0$ y $x_1=1$):
\[
  L_2(x)
    = \frac{(x-0)(x-1)}{(2-0)(2-1)}
    = \frac{x(x-1)}{(2)(1)}
    = \frac{x(x-1)}{2}
\]

\subsection*{Paso 2 \;---\; Construcci\'on de $P(x)$}

\[
  P(x) = \sum_{i=0}^{2} y_i\, L_i(x)
       = 1\cdot L_0(x) + 3\cdot L_1(x) + 0\cdot L_2(x)
\]

Como $y_2 = 0$, el tercer t\'ermino desaparece:
\[
  P(x)
    = 1\cdot\frac{(x-1)(x-2)}{2}
    + 3\cdot\bigl[-x(x-2)\bigr]
\]

\subsection*{Paso 3 \;---\; Expansi\'on y simplificaci\'on}

Expandiendo cada factor:
\[
  (x-1)(x-2) = x^2 - 3x + 2
\]
\[
  P(x)
    = \frac{x^2 - 3x + 2}{2} - 3x^2 + 6x
    = \frac{x^2 - 3x + 2 - 6x^2 + 12x}{2}
    = \frac{-5x^2 + 9x + 2}{2}
\]

\[
  \boxed{P(x) = -\frac{5}{2}\,x^2 + \frac{9}{2}\,x + 1}
\]

\subsection*{Paso 4 \;---\; Verificaci\'on}

\begin{center}
\begin{tabular}{cccc}
  \toprule
  $x_i$ & $P(x_i)$ & $y_i$ & \\
  \midrule
  $0$ & $-\dfrac{5}{2}(0)^2 + \dfrac{9}{2}(0) + 1 = 1$  & $1$ & \checkmark \\[8pt]
  $1$ & $-\dfrac{5}{2}(1)^2 + \dfrac{9}{2}(1) + 1 = 3$  & $3$ & \checkmark \\[8pt]
  $2$ & $-\dfrac{5}{2}(4) + \dfrac{9}{2}(2) + 1 = 0$    & $0$ & \checkmark \\
  \bottomrule
\end{tabular}
\end{center}

El polinomio pasa exactamente por los tres puntos dados.

% ─────────────────────────────────────────────────────────────
\section*{Cota de Error de la Interpolaci\'on}

La f\'ormula general del error de interpolaci\'on de Lagrange con $n+1$ puntos es:
\[
  \bigl|f(x) - P_n(x)\bigr|
  \;\leq\;
  \frac{M_{n+1}}{(n+1)!}\,\prod_{i=0}^{n}|x - x_i|
  \qquad
  M_{n+1} = \max_{\xi\,\in\,[a,b]}\bigl|f^{(n+1)}(\xi)\bigr|
\]

\subsection*{Aplicaci\'on al ejercicio}

Con $n=2$ (tres puntos), el error usa la \textbf{tercera derivada}:
\[
  \bigl|f(x) - P_2(x)\bigr|
  \;\leq\;
  \frac{M_3}{3!}\;\bigl|x(x-1)(x-2)\bigr|
  \qquad
  M_3 = \max_{\xi\,\in\,[0,2]}\bigl|f'''(\xi)\bigr|
\]

\subsection*{Paso 1 \;---\; An\'alisis del t\'ermino de producto $\omega(x)$}

\[
  \omega(x) = x(x-1)(x-2) = x^3 - 3x^2 + 2x
\]

Para hallar su m\'aximo absoluto en $[0,2]$ derivamos e igualamos a cero:
\[
  \omega'(x) = 3x^2 - 6x + 2 = 0
  \;\Longrightarrow\;
  x = 1 \pm \frac{1}{\sqrt{3}}
  \;\Longrightarrow\;
  x_1 \approx 0.4226, \quad x_2 \approx 1.5774
\]

\begin{center}
\begin{tabular}{ccc}
  \toprule
  $x$ & $\omega(x)$ & $|\omega(x)|$ \\
  \midrule
  $0$      & $0$       & $0$ \\
  $0.4226$ & $0.3849$  & $0.3849$ \\
  $1.5774$ & $-0.3849$ & $0.3849$ \\
  $2$      & $0$       & $0$ \\
  \bottomrule
\end{tabular}
\end{center}

Por lo tanto: $\displaystyle\max_{x\,\in\,[0,2]}|\omega(x)| = 0.3849$

\subsection*{Paso 2 \;---\; Cota final}

\[
  \boxed{
    \bigl|f(x) - P_2(x)\bigr|
    \;\leq\;
    \frac{M_3}{6}\times 0.3849
    \;=\;
    0.0642\;M_3
  }
\]

\textbf{Nota:} si $f(x)$ es conocida, se calcula $f'''(x)$ y su m\'aximo en $[0,2]$
para obtener la cota num\'erica exacta. Si $f(x)=e^x$, por ejemplo:
$f'''(x)=e^x$, $M_3=e^2\approx 7.389$, y la cota ser\'ia
$0.0642\times 7.389 \approx 0.474$.

\end{document}
